\section{Construction of the dodecaphase register and its rational phase coordinate}
\label{sec:k-registro}

The generation of the regional coordinates does not exhaust the information
in the APP--TRIT--TPK state. An Archimedean reading determines a value through
compatible cylinders; another reading preserves the opposition of the sheets,
the orientation of the route, and the record of the incidences traversed. The
dodecaphase register belongs to this second reading before entering the closure
of alpha. It is therefore not defined by subtracting a physical constant from
a combination of previously known values. It is reconstructed from a signed
coordinate of the same arithmetic state with path memory that produces the
regional coordinates.

The distinction is necessary even when two objects have twelve coordinates.
A twelve-block word from the catalog, a list of twelve oriented events, a
signed integral vector, and a base-one-thousand dodecaphase register are
distinct objects. Comparing them requires the mappings that relate them.
Equality of length does not supply those mappings, and the loss of signs is
not a mere simplification of notation.

This chapter maintains the causal cut fixed in the preceding chapters:
the nine APP positions of the horizontal and vertical axes produce the
arithmetic sheets; TRIT types regime and orientation; TPK selects,
transports, updates, and prolongs, preserving remainder, quotient, carry,
sheet, route, boundary, and memory. The coordinates used here are obtained
from the joint discrete structure of the continuum. The chain
\(w_6\to w_{12}\to w_{18}\to w_{24}\to w_{30}\to R_{36}\to G_9\)
is not replaced by a finite calendar. Nor is the order
\(U_t\to\Gamma_9\to\mathcal K_{\mathrm{ph}}\) reversed: the last mapping is
a reduced-memory reading, not the generator of the signed record.

% Proposed addition at the beginning of registro_k.tex, after its introduction.
% Preserves all current definitions, formulas, and proofs.
% Provenance: registro_k.tex, §§k-hadamard/k-transversal/k-kappas;
% k_reversibilidad.tex; incidencia_articulacion.tex, inc-ampl-cuaterna.
\paragraph{Reversible coordinatization and incidence determination.}
\label{ampl-alfa-k-intro-reversible}
The result organizing this chapter is the equivalence among three
coordinatizations of a record previously produced by
APP--TRIT--TPK. The Hadamard transformation relates the signed record
to the integral record; the difference extractor relates
the latter to its transverse variations and its uniform component:
\[
 U\underset{\mathcal H_{12}}{\overset{\mathcal H_{12}/4}{\longleftrightarrow}}K
 \overset{(D_3,D_4,Q)}{\longleftrightarrow}
 (D_3K,D_4K,Q(K)).
\]
The first equivalence is established on the subset of the integrality
sublattice $\mathcal L_H$ that constitutes the domain of this chart; the
second, on the image of the extractor. Once the base $1000$, the $12$
positions, and the phase origin are fixed, the evaluation
$K\mapsto\kappa_{\mathrm{per}}$ adds a fourth reversible
coordinatization: multiplication by $1000^{12}-1$ followed by Euclidean
division recovers the complete record. The proofs of
\S\ref{subsec:k-hadamard}, \S\ref{subsec:k-transversal}, and
\S\ref{subsec:k-kappas} establish these equivalences on their respective
domains.

Two compositions subsequently act on that record. The positional
composition combines its components with the closure, propagation, and
self-scaling coordinates, and its normalization determines $\alpha_A$. The
incidence composition applies the cubic threshold and obtains
\[
 K\longmapsto B_K=\{j:K_j\geq729\}=\{4,6,9,10\}.
\]
The complete record makes it possible to reconstruct both compositions.
The support $B_K$ specifically preserves membership in the class selected
by the threshold. The identity $B_K=H_e\cap P_\pi$ in
\eqref{inc-ampl-cuaterna} links this integral extraction to the incidence
of the regional codes. The negative support of its balance with $K$
completes the flag $B_K\subset H^-_\alpha$ developed in the exceptional
prolongation. Thus, the function of $K$ is defined through an invertible
transformation and precise subsequent mappings; the additional depth and
transport information remains in the state from which the record arises.


\subsection{Cyclic extension and memory grading}
\label{subsec:k-calendario}

The observable calendar contains one hundred eight positions, organized into
twelve nonadic windows. With positive indices, let
\[
 E_m=\{9(m-1)+1,\ldots,9m\},\qquad 1\leq m\leq12.
\]
The ordered support \(E_1\sqcup\cdots\sqcup E_{12}\) preserves position and
window membership. The group \(C_{108}=\mathbb Z/108\mathbb Z\), by
contrast, preserves the cyclic translation law. For the representative
\(0\leq\widetilde\theta<108\), the coordinates
\[
 \beta(\theta)=\left[\left\lfloor\frac{\widetilde\theta}{9}\right\rfloor\right]_{12},
 \qquad r(\theta)=1+(\widetilde\theta\bmod9)
\]
separate the dodecaphase sector and the interior nonadic position. Neither
of them by itself preserves the total number of turns.

\begin{proposition}[The carry of the calendar]
\label{prop:k-extension}
The mappings
\[
 \iota:C_{12}\longrightarrow C_{108},\quad[b]\longmapsto[9b],
 \qquad p:C_{108}\longrightarrow C_9,\quad[\theta]\longmapsto[\theta]_9
\]
form a nonsplit exact sequence
\[
 0\longrightarrow C_{12}\xrightarrow{\iota}C_{108}
 \xrightarrow{p}C_9\longrightarrow0.
\]
For the set-theoretic section that chooses the representatives
\(0,\ldots,8\), the additivity defect is the carry
\[
 c(r,r')=\left[\left\lfloor\frac{\widetilde r+\widetilde r'}9\right\rfloor\right]_{12}.
\]
\end{proposition}

\begin{proof}
The kernel of \(p\) consists of the multiples of nine and coincides with the
image of \(\iota\). The reduction is surjective, and multiplication by nine
is injective on the indicated domain. If the sequence split, the central
group would be isomorphic to \(C_{12}\times C_9\), whose exponent is
\(\operatorname{lcm}(12,9)=36\); the exponent of \(C_{108}\) is 108. The
contradiction proves that the sequence is nonsplit. Finally, Euclidean
division of \(\widetilde r+\widetilde r'\) by nine gives the displayed carry
and verifies the identity of the section.
\end{proof}

The result explains a difference that will reappear in alpha: returning to
the phase is not equivalent to returning to the state. If
\(q_{\mathrm{mem}}\) counts the nonadic turns, the calendar preserves only
\(\beta=[q_{\mathrm{mem}}]_{12}\). After one hundred eight steps,
\(\beta\) returns, but \(q_{\mathrm{mem}}\) increases by twelve. The record
also preserves the cylinder, the incidence ledger, and the vertical memory.
The closure of a finite chart does not turn that history into a memoryless
cycle.

\subsection{The oriented record and the integral \texorpdfstring{\(1+5+6\)}{1+5+6} chart}
\label{subsec:k-registro-integral}

For a history produced by TPK, let \(\tau_m\) be the subroute on
\(E_m\), \(\sigma_m\) its orientation, \(\kappa_m\) the boundary carry,
and \(\Lambda_m\) the local incidence ledger. The record is
\[
 \mathcal R_{12}(x)=
 \bigl((E_m,\tau_m,\sigma_m,\kappa_m,\Lambda_m)\bigr)_{m=1}^{12}.
\]
Its domain is the space of arithmetic histories with path memory, not the set
of phase labels. The projection to the calendar forgets precisely some of
the data that the signed reader requires.

An initial chart of the record can be written through oriented events. The
construction developed in this article fixes the set of event codes
\(\{2,4,6,8\}\) and the sectorial orientation
\[
 \Sigma_{12}=(+,+,+,-,-,-,+,+,+,-,-,-).
\]
If \(d_t\) is the direction code at step \(t\), one obtains
\[
 e_i=\Sigma_{12}(i+1)
 \#\{t\in E_{i+1}:d_t\in\{2,4,6,8\}\},
 \qquad0\leq i<12.
\]
The formula keeps the nonnegative event census and the orientation separate.
It does not identify a cardinal direction with a TRIT sign. Nor is the vector
\(e\), whose components are signed local censuses, yet the vector \(U\) that
will appear below.

The opposite positions determine, for \(0\leq i<6\),
\[
 p_i=e_i+e_{i+6},\qquad a_i=e_i-e_{i+6}.
\]
With
\[
 s_C=\sum_{i=0}^{5}p_i,\qquad M_i=p_i-p_5\quad(0\leq i<5),
\]
define the integral chart
\[
 \mathcal L_{12}(e)
 =(s_C,M_0,\ldots,M_4,a_0,\ldots,a_5).
\]
The \(1+5+6\) decomposition is not an arbitrary distribution of coordinates:
one coordinate preserves the charge, five compare the pairs relative to a
reference pair, and six preserve the opposition of the half-coronas.

\begin{lemma}[Inversion of the integral chart]
\label{lem:k-carta-integral}
An integral tuple \((s_C,M_0,\ldots,M_4,a_0,\ldots,a_5)\) belongs to the
image of \(\mathcal L_{12}\) if and only if
\[
 s_C-\sum_{i=0}^{4}M_i\equiv0\pmod6,
 \qquad p_i\equiv a_i\pmod2\quad(0\leq i<6),
\]
where
\[
 p_5=\frac{s_C-\sum_{i=0}^{4}M_i}{6},\qquad p_i=p_5+M_i\quad(i<5).
\]
In that case its preimage is unique and is given by
\[
 e_i=\frac{p_i+a_i}{2},\qquad e_{i+6}=\frac{p_i-a_i}{2}.
\]
\end{lemma}

\begin{proof}
The sum of the five margins is \(s_C-6p_5\), which yields the division by
six. Once the \(p_i\) have been reconstructed, each pair of equations
\(p_i=e_i+e_{i+6}\), \(a_i=e_i-e_{i+6}\) has the stated solution. This
solution is integral exactly under the parity congruence. All operations are
inverted, so no residual freedom remains.
\end{proof}

This reversibility shows what information is preserved; it does not authorize
the suppression of the rest of the ledger. The complete record has a larger
domain than this census chart. The extraction of the channels that feed
\(U\) uses that complete record and is not deduced from the list \((e_i)\)
through an identification of names.

\subsection{Positional and incidence readings of the terminal state}
\label{subsec:k-terminal}

Denote by \(X_{\mathrm{term}}\) the state produced by the composition of
selection, transport, and updating up to the terminal cut. The operator
expression is
\[
 X_{\mathrm{term}}=
 (\mathcal U_{108}\circ\cdots\circ\mathcal U_1)(X_{\mathrm{seed}}),
 \qquad \mathcal U_t=\mathsf{Upd}_t\circ\mathsf{Tra}_t\circ\mathsf{Sel}_t.
\]
Its two relevant readings are
\[
 X_{\mathrm{term}}
 \xrightarrow{(\pi_{\mathrm{term}},\mathcal S_{12})}
 (B_{\mathrm{term}},U).
\]
The matrix \(B_{\mathrm{term}}\) retains the visible ternary reading; the
vector \(U\) preserves another coordinate of the same history. No arrow
\(B_{\mathrm{term}}\to U\) is inserted between them.

The terminal selector developed in this article operates on catalogs of
sizes 196, 197, and 196. The local signatures reduce their
\(196\cdot197\cdot196=7\,567\,952\) triples to 331, and the column margin
\(q_{\partial}=(1,2,2,1,2,0)\) leaves two matrices:
\[
 B_0=\begin{pmatrix}0&2&0&2&2&0\\0&0&1&2&2&0\\1&0&1&0&1&0\end{pmatrix},
 \qquad
 B_1=\begin{pmatrix}0&2&1&0&2&0\\0&0&1&2&2&0\\1&0&0&2&1&0\end{pmatrix}.
\]
The row orientation is \(\sigma=(1,1,-1)=(1,1,2)\) in
\(\mathbb F_3^3\). The row charges of the candidates are \((0,2,0)\) and
\((2,2,1)\). Only the second belongs to
\(\langle\sigma\rangle=\{(0,0,0),(1,1,2),(2,2,1)\}\); therefore
\(B_{\mathrm{term}}=B_1\). This verification proves the final step of the
selector from the stated census. It does not turn the column margin into a
consequence of the row charge, nor does it replace the preceding enumeration
with inspection of those two matrices.

The signed reader is the composition
\begin{equation}
 \mathcal S_{12}
 =\Pi_H\circ\mathcal E_{108}^{90,120}\circ\mathcal R_{12}:
 \mathcal X_{\mathrm{TPK}}^{\mathrm{term,mem}}
 \longrightarrow\mathbb Z^{12}.
 \label{eq:k-lector-firmado}
\end{equation}
Its evaluation is developed below on the complete enriched state, without
introducing \(K\), \(\alpha\), or conventional metrology as
selectors. The raw CT108 projection, which forgets the signed coordinate, is
not used as a substitute domain for this reader.

% Autonomous typing of the terminal dodecaphase chart.
\subsection{Signed dodecaphase chart of the full TPK state}
\label{subsec:registro-evaluacion-incidencial}

The dodecaphase chart and a possible factorization through the raw
108-step projection (CT108) are objects of different types. Let
\(X\in\mathsf{TPKFullState}\) be a complete enriched state, and let
\(\rho_{108}(X)\in\mathsf{CT108RawProjection}\) be the object obtained by
forgetting orientation, sheet, carry, boundary signature, memory, and the
dodecaphase ledger. In the canonical coordinates of the full state, write
\[
 X=(X_{\mathrm{APP}},X_{\mathrm{TRIT}},X_{\mathrm{route}},X_{\mathrm{mem}},
       C_{\mathrm{term}},U_{12}^{\mathrm{sgn}},
       X_{\mathrm{boundary}},\ldots).
\]
The active reader is, by definition, the projection of that already generated
and preserved coordinate:
\[
 \mathcal S_{12}:=\operatorname{pr}_{U_{12}^{\mathrm{sgn}}}:
 \mathsf{TPKFullState}\longrightarrow\mathsf{SignedDodecaphaseState}.
\]
On the compatible sublattice of the full state, this projection coincides
with the composite reader
\begin{equation}
 \mathcal S_{12}\equiv R_{\mathrm{sgn}}
 :=\Pi_H\circ\mathcal E_{108}^{90,120}\circ\mathcal R_{12}.
 \label{eq:registro-identidad-s12-rsgn}
\end{equation}
Therefore, \(\mathcal S_{12}\) reads a preserved coordinate of the same state;
it receives neither a conventional constant nor a metrological table. Its
domain is not replaced by \(\mathsf{CT108RawProjection}\), since that
replacement would additionally require a factorization
\(\mathcal S_{12}=\Xi\circ\rho_{108}\). Such a factorization is neither
presupposed nor needed here; it would require proof that the signed reading is
constant on the fibers of \(\rho_{108}\). The internal theorem operates
directly on the full TPK state.

To fix the content of a possible incidence realization, let \(\mathscr L\) be
a ledger of visits and oriented traces belonging to \(X\). Each visit
preserves the route, APP position, sheet, instant, multiplicity, and boundary
incidence. The arithmetic evaluation produces the integer and its
decomposition
\[
 n_v=r_v+9q_v,\qquad 1\leq r_v\leq9.
\]
The integer \(n_v\) is calculated from the sum or product sheet; its remainder
is not received as an independent value. The distinction between corona
incidence and interior incidence remains explicit for visits of remainder
\(9\).

In the temporal window
\(E_m=\{9(m-1)+1,\ldots,9m\}\), the domains are distinguished
\[
 \mathscr V_m=\{v:t(v)\in E_m\},\qquad
 \mathscr T_m=\{\theta:\operatorname{origin}(\theta)\in E_m\}.
\]
The presentation by counts considers
\begin{align}
 A_m&=\sum_{v\in\mathscr V_m}w(v)
             \mathbf1_{\{1,3,5,7\}}(r(v)),\\
 C_m&=-\sum_{\theta\in\mathscr T_m}u(\theta)\sigma(\theta)
      =\sum_{j\geq0}\bigl(\nu^-_{120,m}(j)-\nu^+_{120,m}(j)\bigr),\\
 V_m&=\sum_{v\in\mathscr V_m}w(v)
              \mathbf1_{\{9\}}(r(v))\epsilon_\partial(v).
\end{align}
Here \(w(v)\) is the incidence multiplicity, and
\(\epsilon_\partial(v)=1\) on the corona, \(-1\) in the interior.
The traces counted by \(\nu^\pm\) must be enumerated by the corresponding
ledger rule. The integral sum requires finite support or an alternative
construction that determines its meaning. The number of chronological
incidences and the reduced length \(120(1+3j)\) have different types; their
relation requires the reader's unit of length.

Define \([x]_{10}\in\{0,\ldots,9\}\) as the Euclidean representative and
preserve the quotients
\[
 x=[x]_{10}+10\lfloor x/10\rfloor.
\]
The incidence block is
\begin{equation}
 z_m=100[A_m]_{10}+10[C_m]_{10}+[V_m]_{10}.
 \label{eq:registro-bloque-incidencial}
\end{equation}
By composition one obtains the mapping
\begin{equation}
 \mathcal E_{\rm cnt}(\mathscr L)=
 \left((z_m-z_{m+3})_{m=1}^{12},
       (z_m-z_{m+4})_{m=1}^{12},\sum_{m=1}^{12}z_m\right).
\label{eq:registro-composicion-incidencial}
\end{equation}
The indices are read cyclically modulo \(12\). In this subsection,
\((D_jz)_m=z_m-z_{m+j}\), for \(j=3,4\), and
\(Q(z)=\sum_{m=1}^{12}z_m\); therefore, the preceding composition is
\((D_3z,D_4z,Q(z))\). This convention is preserved in
\S\ref{subsec:k-transversal}, where the complete integral reconstruction is
proved. This formula produces its coordinates from the counts. It uses
neither \(K\), \(U\), \(\alpha\), nor the expected transverse coordinates as
arguments.

\subsubsection{Terminal representation and exact closure of its charts}

The canonical terminal state is
\[
 X_{\mathrm{term}}
 =\bigl(\mathcal U_{108}\circ\cdots\circ\mathcal U_1\bigr)
   (X_{\mathrm{seed}}),
 \qquad
 \mathcal U_t=\mathsf{Upd}_t\circ\mathsf{Tra}_t\circ\mathsf{Sel}_t,
\]
where each \(\mathcal U_t\) acts on the full TPK state and not on its raw
projection. Among the coordinates produced and preserved before the signed
representation is
\[
 C_{\mathrm{term}}=\operatorname{pr}_{C}(X_{\mathrm{term}})
 =(b^{90},b^{120},Q_{\mathrm{TPK}}),
\]
with
\begin{align*}
 b^{90}={}&(-495,-116,-684,108,601,-90,
              -173,-88,313,560,-397,461),\\
 b^{120}={}&(-425,-281,-481,671,-255,30,
               475,-543,680,251,6,-128),\\
 Q_{\mathrm{TPK}}={}&6263.
\end{align*}
No subsequent table is introduced here: \(C_{\mathrm{term}}\) is the
transverse coordinate that \(\mathcal E_{108}^{90,120}\) produces within the
evolution of the full state. On the compatible sublattice of TPK states, the
chart satisfies the operator identity
\begin{equation}
 \mathcal S_{12}=\operatorname{pr}_{U}
 =\Pi_H\circ\operatorname{pr}_{C}.
 \label{eq:registro-identidad-s12-pih}
\end{equation}
The identity is formulated on the full TPK state and does not presuppose the
descent of \(\mathcal S_{12}\) through \(\rho_{108}\).

The subsequent evaluation from the twenty-five expressed fields is displayed
without a gap. For \(1\le r\le3\), let
\[
 B_r=\sum_{j=0}^{3}b^{120}_{r+3j},\qquad
 d_{r,j}=b^{90}_{r+3j},
\]
and
\[
 s_1=\frac{Q_{\mathrm{TPK}}+2B_1+B_2}{3},\qquad
 s_2=s_1-B_1,\qquad s_3=s_2-B_2.
\]
Integral substitution of the preceding twenty-five fields gives
\[
 (B_1,B_2,B_3)=(972,-1073,101),\qquad
 (s_1,s_2,s_3)=(2378,1406,2479).
\]
Applying
\[
 U^{(r)}=(s_r,d_{r,1}-d_{r,3},d_{r,0}+d_{r,2},d_{r,0}-d_{r,2})
\]
one obtains, component by component,
\begin{align*}
 U^{(1)}&=(2378,-452,-668,-322),\\
 U^{(2)}&=(1406,998,-204,-28),\\
 U^{(3)}&=(2479,-551,-371,-997).
\end{align*}
Columnwise interleaving therefore evaluates the signed chart:
\begin{equation}
 U_{\mathrm{term}}:=\mathcal S_{12}(X_{\mathrm{term}})
 =(2378,1406,2479,-452,998,-551,-668,-204,-371,-322,-28,-997).
 \label{eq:registro-uterm-desde-estado-pleno}
\end{equation}

\begin{theorem}[Subsequent exact closure among \texorpdfstring{\(U\)}{U}, \texorpdfstring{\(K\)}{K}, and the channels]
\label{thm:registro-evaluacion-terminal-recuperada}
Once the signed chart \eqref{eq:registro-uterm-desde-estado-pleno} has been
expressed by projection, the orbital Hadamard transform expresses, at a
subsequent stage,
\[
 K=(234,543,140,729,659,824,621,058,914,794,146,601),
\]
and the transverse extractor \((D_3,D_4,Q)\) produces
\begin{align}
 b^{90}={}&(-495,-116,-684,108,601,-90,
             -173,-88,313,560,-397,461),\notag\\
 b^{120}={}&(-425,-281,-481,671,-255,30,
              475,-543,680,251,6,-128),\notag\\
 Q_{\mathrm{TPK}}={}&6263.
 \label{eq:registro-canales-terminales-producidos}
\end{align}
The composition is an exact bijection among the compatible sublattice of signed
states, the seal \(K\), and its transverse coordinate
\(C_{\mathrm{term}}=(b^{90},b^{120},Q_{\mathrm{TPK}})\).
In particular,
\begin{equation}
 C_{\mathrm{term}}=(D_3K,D_4K,Q(K)).
 \label{eq:registro-cterm-desde-xterm}
\end{equation}
\end{theorem}

\begin{proof}
Grouping \(U_{\mathrm{term}}\) into its three step-three orbits yields
\[
 \begin{aligned}
  u^{(0)}&=(2378,-452,-668,-322),&
  \tfrac14H_4u^{(0)}&=(234,729,621,794),\\
  u^{(1)}&=(1406,998,-204,-28),&
  \tfrac14H_4u^{(1)}&=(543,659,58,146),\\
  u^{(2)}&=(2479,-551,-371,-997),&
  \tfrac14H_4u^{(2)}&=(140,824,914,601).
 \end{aligned}
\]
The divisions are exact in \(\mathbb Z\).
Columnwise interleaving expresses \(K\). Finally, the twelve cyclic
step-three and step-four subtractions and the total sum give, component by
component, \eqref{eq:registro-canales-terminales-producidos}. Proposition
\ref{prop:k-transversal} proves the integral inverse, and
Theorem~\ref{thm:k-naturalidad-lector-firmado} proves naturality under
refinement. No operation receives \(\alpha\), CODATA, a metrological table,
or a conventional constant.
\end{proof}

\begin{remark}[Complete operator form and focal executable control]
The terminal arrow is not postulated from a table: it is the typed
composition
\[
 X_{\mathrm{term}}
 \xrightarrow{\mathcal R_{12}}\mathscr L(X_{\mathrm{term}})
 \xrightarrow{\mathcal E_{108}^{90,120}}C_{\mathrm{term}}
 \xrightarrow{\Pi_H}U_{\mathrm{term}}
 \xrightarrow{\mathcal H_{12}^{-1}}K .
\]
The preceding definitions traverse the twelve windows that partition the one
hundred eight steps and express each channel coordinate as a finite sum of its
visits and traces, with orientation, sheet, carry, boundary, and memory. This
is the mathematical realization of the production of
\(C_{\mathrm{term}}\); a row-by-row tabular expansion would be another
serialization of the same sum, not an additional premise of the theorem. The
focal executable exactly reevaluates the finite chart
\(C_{\mathrm{term}}\leftrightarrow U_{\mathrm{term}}\leftrightarrow K\)
as a supplementary reproducible control. The proof of the continuum
hypothesis remains formulated over APP--TRIT--TPK, the full enriched state,
its correlated representations, and the joint discrete structure.
\end{remark}

Neither \(\mathcal S_{12}\), nor \(U\), nor \(K\), nor \(\alpha\) are
external inputs or conventional selectors of the discrete construction of the
continuum. They are typed coordinates and representations of the same
APP--TRIT--TPK state. The coordinate \(\alpha_{\mathrm{HMT}}\) is a derived
and correlated output of that single generation; its representation order does
not break the continuity of the object or weaken its derivation.

The incidence formula \eqref{eq:registro-composicion-incidencial} describes an
additional realization when a complete ledger
\(\mathscr L(X_{\mathrm{term}})\) is unfolded. The preceding theorem does not
claim that the signed coordinate can be reconstructed after that ledger has
been forgotten through \(\rho_{108}\); nor does it require such a claim. This
separation avoids confusing the internal closure of the enriched state with
an inverse problem posed on a projection that has removed precisely the
orientation and memory information.

The mapping \(\Pi_H\), executed within the chart, produces from
\(C_{\mathrm{term}}\) the signed coordinate and its three step-three orbits:
\begin{align*}
 U_{\mathrm{term}}^{(1)}&=\Pi_H^{(1)}(C_{\mathrm{term}})
                         =(2378,-452,-668,-322),\\
 U_{\mathrm{term}}^{(2)}&=\Pi_H^{(2)}(C_{\mathrm{term}})
                         =(1406,998,-204,-28),\\
 U_{\mathrm{term}}^{(3)}&=\Pi_H^{(3)}(C_{\mathrm{term}})
                         =(2479,-551,-371,-997).
\end{align*}
Let \(\mathcal P_{\mathrm{orb}}:\mathbb Z^{12}\to(\mathbb Z^4)^3\) be the
orbital regrouping
\[
 \mathcal P_{\mathrm{orb}}(U)
 =\bigl((U_1,U_4,U_7,U_{10}),
        (U_2,U_5,U_8,U_{11}),
        (U_3,U_6,U_9,U_{12})\bigr),
\]
and let \(\operatorname{Int}_{3,4}=\mathcal P_{\mathrm{orb}}^{-1}\) be the
columnwise interleaving of three tetrads. We also define
\[
 \operatorname{Dig}_{10}^{(3)}(z)
 =\left(\left\lfloor\frac z{100}\right\rfloor,
         \left\lfloor\frac z{10}\right\rfloor\bmod10,
         z\bmod10\right),\qquad0\le z\le999.
\]
The terminal modular chart is the effective mapping on the compatible channel
state
\begin{equation}
 \mathcal G_{\mathrm{term}}(C)
 :=\left[(\operatorname{Dig}_{10}^{(3)})^{12}
 \circ\operatorname{Int}_{3,4}
 \circ\left(\tfrac14H_4\oplus\tfrac14H_4
                    \oplus\tfrac14H_4\right)
  \circ\mathcal P_{\mathrm{orb}}\right]\!\left(\Pi_H(C)\right).
 \label{eq:registro-generador-modular-terminal}
\end{equation}
Its domain is the subset of the integrality sublattice of compatible
coordinates in \(\mathbb Z^{25}\) on which \(\Pi_H\) is integral, the three
orbital images under \(H_4\) are divisible by four, and the twelve resulting
coordinates belong to \(\{0,\ldots,999\}\). It receives neither
\(U_{\mathrm{term}}\), \(K\), \(\alpha\), a decimal table, nor a metrological
value; the interleaving, the phase origin, and the orientation are part of the
type of the channel coordinate.

The verification coordinate is serialized in
\path{metadata/semilla-registro-terminal.json}. That file contains
\(C_{\mathrm{term}}\), the phase origin, and the orientation; it contains
neither \(U_{\mathrm{term}}\), \(K\), nor the twelve residual rows. The focal
program applies \(\mathcal G_{\mathrm{term}}\), expresses
\(U_{\mathrm{term}}\), \(K\), and the residues, and returns through
\((D_3,D_4,Q)\) to the same \(C_{\mathrm{term}}\). Starting there permits an
exact focal verification of the subsequent finite chart. It does not redefine
the domain of the signed chart
\eqref{eq:registro-uterm-desde-estado-pleno}, does not replace the preceding
proof, and does not attribute to the auxiliary the production of the
coordinate it receives.

\begin{proposition}[Modular calculation of the twelve terminal rows]
\label{prop:registro-generacion-modular-terminal}
The evaluation of \(\mathcal G_{\mathrm{term}}\) on the produced coordinate
\(C_{\mathrm{term}}\) first obtains \(U_{\mathrm{term}}\) and then the phase
tetrads
\[
 (234,729,621,794),\qquad(543,659,058,146),\qquad
 (140,824,914,601),
\]
and, after interleaving, produces exactly the twelve rows of the following
table. It is the algebraic verification subsequent to
Theorem~\ref{thm:registro-evaluacion-terminal-recuperada}; it does not replace
the representation of the signed chart of the full TPK state.
\end{proposition}

\begin{proof}
The formulas for \(\Pi_H\) give the three signed blocks written above. The
subsequent product by \(H_4/4\) gives, without rounding,
\begin{align*}
 \tfrac14H_4(2378,-452,-668,-322)^T&=(234,729,621,794)^T,\\
 \tfrac14H_4(1406,998,-204,-28)^T&=(543,659,58,146)^T,\\
 \tfrac14H_4(2479,-551,-371,-997)^T&=(140,824,914,601)^T.
\end{align*}
The twelve divisions are exact in \(\mathbb Z\). Columnwise interleaving gives
\[
 \begin{split}
 (&234,543,140,729,659,824,\\
  &621,058,914,794,146,601),
 \end{split}
\]
and \(\operatorname{Dig}_{10}^{(3)}\) uniquely provides its three decimal
residues. The auxiliary verifier included in the package executes these same
operations; its negative controls alter the channel state or its phase origin.
\end{proof}

Denote by
\[
 \mathscr F_{\mathrm{term}}^{(10)}
 :=\bigl((a_m,c_m,v_m)\bigr)_{m=1}^{12}
\]
the digit chart expressed by \(\operatorname{Dig}_{10}^{(3)}(K_m)\).
When a complete incidence realization is available,
\((a_m,\allowbreak c_m,\allowbreak v_m)=([A_m]_{10},\allowbreak[C_m]_{10},\allowbreak[V_m]_{10})\).

This equality is not
postulated in order to reconstruct the ledger after it has been forgotten.
Each row preserves the window index and is therefore not a bag of
interchangeable digits.

\begin{center}
\begin{tabular}{c@{\qquad}ccc@{\qquad}c}
\toprule
\(m\)&\(a_m\)&\(c_m\)&\(v_m\)&\(K_m\)\\
\midrule
1&2&3&4&234\\
2&5&4&3&543\\
3&1&4&0&140\\
4&7&2&9&729\\
5&6&5&9&659\\
6&8&2&4&824\\
7&6&2&1&621\\
8&0&5&8&058\\
9&9&1&4&914\\
10&7&9&4&794\\
11&1&4&6&146\\
12&6&0&1&601\\
\bottomrule
\end{tabular}
\end{center}
Consequently,
\begin{equation}
 z^{\mathrm{term}}
 =(234,543,140,729,659,824,621,058,914,794,146,601).
 \label{eq:registro-terminal-residuos}
\end{equation}

\begin{proposition}[Terminal return of the computed collection to the channels]
\label{prop:registro-terminal-coordenadas}
The decimal recomposition of \(\mathscr F_{\mathrm{term}}^{(10)}\), followed
by the transverse extractor, produces, without consulting \(\alpha\) or a
metrological table,
\begin{align*}
 D_3z^{\mathrm{term}}={}&(-495,-116,-684,108,601,-90,
 &\hspace{18mm}-173,-88,313,560,-397,461),\\
 D_4z^{\mathrm{term}}={}&(-425,-281,-481,671,-255,30,
 &\hspace{18mm}475,-543,680,251,6,-128),\\
 Q(z^{\mathrm{term}})={}&6263.
\end{align*}
These are exactly the twenty-five coordinates
\((b^{90},b^{120},Q_{\mathrm{TPK}})\) used by the harmonic reader.
\end{proposition}

\begin{proof}
The decimal formation of each row gives
\eqref{eq:registro-terminal-residuos}. Cyclically subtracting the entry
\(m+3\) and the entry \(m+4\) produces, component by component, the two
vectors displayed; the sum of the twelve entries is \(6263\). The calculation
is repeated in the package's autonomous verifier. The terminal collection is a
computed representation of the channel state. The transverse evaluation returns
exactly to the seed \(C_{\mathrm{term}}\); the harmonic reading then recovers
\(U_{\mathrm{term}}\), and the inverse Hadamard transform recovers the same
\(K\). This proves the agreement of both coordinatizations without using
alpha as a selector.
\end{proof}

The chart \(\mathscr F_{\mathrm{term}}^{(10)}\) is sufficient and necessary
for this transverse evaluation. The table is a decimal certificate of the
expressed output, not a visit-by-visit reedition of the oriented ledger. The
enriched state additionally preserves Euclidean quotients, visits, traces,
and provenance; those fields are not declared recoverable from the thirty-six
residues, nor are they necessary for the cardinal theorem of this article.

\begin{proposition}[Sufficient and necessary arithmetic information]
\label{prop:registro-36-residuos}
Let \(N,N'\in\mathbb Z^{12\times3}\), ordered by the counts
\((A,C,V)\). For the mapping defined by
\eqref{eq:registro-bloque-incidencial}--\eqref{eq:registro-composicion-incidencial},
\[
 \mathcal E_{\rm cnt}(N)=\mathcal E_{\rm cnt}(N')
 \quad\Longleftrightarrow\quad
 N-N'\in10\mathbb Z^{36}.
\]
In particular, the \(36\) sectorial residues exactly determine the output of
\(25\) coordinates. The preservation of the quotients and provenance
constitutes additional information in the record.
\end{proposition}

\begin{proof}
Equality modulo \(10\) produces the same blocks and, therefore, the same
transverse coordinates. Conversely, if the transverse coordinates coincide,
\(D_3(z-z')=D_4(z-z')=0\). The steps \(3\) and \(4\) generate the cycle of
length \(12\), so \(z-z'\) is constant. Equality of the total charges makes
that constant vanish. Finally, the decimal representation of an integer
between \(0\) and \(999\) has three unique digits. The three residues of each
window are thus recovered. The assertion is an algebraic identity for every
\(N,N'\), not an inference obtained from finite tests.
\end{proof}

\subsection{Conjugation of the counts and boundary terms}
\label{subsec:registro-conjugacion-incidencial}

Let \(\rho(m)=13-m\). Consider an incidence conjugation that reverses the
order of the windows, preserves the arithmetic class and the boundary class
of the visits, and interchanges the trace channels. Then
\[
 (A_m^\dagger,C_m^\dagger,V_m^\dagger)
 =(A_{\rho(m)},-C_{\rho(m)},V_{\rho(m)}).
\]
If \(c_m=[C_m]_{10}\), decimal reduction produces the identity
\begin{equation}
 z_m^\dagger=z_{\rho(m)}
       +100\mathbf1_{\{c_{\rho(m)}\ne0\}}-20c_{\rho(m)}.
 \label{eq:registro-conjugacion-decimal}
\end{equation}
Indeed, \([-C]_{10}=0\) when \([C]_{10}=0\), and
\([-C]_{10}=10-[C]_{10}\) otherwise. The formula preserves exactly this
difference. Negating the channel is not equivalent to negating the complete
decimal block.

Applying this identity to a concrete TPK involution requires verifying its
action on the incidences. For an occupancy reader evaluated before the
advance, the reversal of a route
\(\gamma=(x_0,\ldots,x_n)\) satisfies
\[
 \sum_{k=0}^{n-1}f(x_{n-k})
 =\sum_{k=0}^{n-1}f(x_k)+f(x_n)-f(x_0).
\]
The endpoint terms are incorporated before reduction modulo \(10\). They
vanish on closed routes; in open windows they may alter the residues. This
correction, already present in the bidirectional development of transport,
specifies what the linkage between the route and the counts must preserve.

\begin{proposition}[Necessity of the nonperiodic information in the record]
Suppose that the observable of a fixed collection of routes satisfies
\(o(t+54)=o(t)\), that multiplicity is preserved under that return, and that
the same local reader acts on the windows of \(9\) instants. Then its counts
and blocks satisfy
\[
 (A,C,V)_{m+6}=(A,C,V)_m,\qquad z_{m+6}=z_m.
\]
\end{proposition}

\begin{proof}
The translation \(t\mapsto t+54\) carries \(E_m\) bijectively to
\(E_{m+6}\) and preserves each counted incidence and its multiplicity. The
equality is transmitted to the sum and to decimal reduction.
\end{proof}

Therefore, a presentation with \(12\) distinct blocks requires the reader to
preserve some information that does not factor through that observable of
period \(54\): incidence selection, effective orientation, endpoints,
multiplicity, or memory. The proposition delimits an inadmissible reduction
of the state; it does not assert that the complete dynamics has period \(54\),
nor does it by itself determine its terminal reader.


Theorem~\ref{thm:registro-evaluacion-terminal-recuperada} expresses the
signed chart of the full TPK state, applies the orbital Hadamard transform,
and obtains through the transverse extractor
\begin{align*}
 b^{90}={}&(-495,-116,-684,108,601,-90,\\
            &\hspace{19mm}-173,-88,313,560,-397,461),\\
 b^{120}={}&(-425,-281,-481,671,-255,30,\\
             &\hspace{19mm}475,-543,680,251,6,-128),\\
 Q_{\mathrm{TPK}}={}&6263.
\end{align*}
These quantities are internal outputs of the terminal chart, not parameters
fitted through alpha. Formulas \eqref{eq:k-sumas-armonicas}--
\eqref{eq:k-bloques-firmados} recover from them
\begin{equation}
 U_{\mathrm{term}}
 =\mathcal S_{12}(X_{\mathrm{term}})
 =(2378,1406,2479,-452,998,-551,
   -668,-204,-371,-322,-28,-997).
 \label{eq:k-u-terminal-directo}
\end{equation}
Proposition~\ref{prop:registro-terminal-coordenadas} also verifies the
transverse return to the same twenty-five channels. The arithmetic preserves
the distinction among the terminal matrix, the incidence collection, the signed
vector, and the dodecaphase register.

\subsection{Construction of the signed vector from the channel record}
\label{subsec:k-pi-h}

The harmonic reader \(\Pi_H\) can be written without using either \(K\) or
alpha. To avoid confusing its orbital sums with the alpha blocks, denote them
by \(s_1,s_2,s_3\). Let
\[
 B_r=\sum_{j=0}^{3}b^{120}_{r+3j},\qquad
 d_{r,j}=b^{90}_{r+3j},\qquad1\leq r\leq3,
\]
and
\begin{equation}
 s_1=\frac{Q_{\mathrm{TPK}}+2B_1+B_2}{3},\qquad
 s_2=s_1-B_1,\qquad s_3=s_2-B_2.
 \label{eq:k-sumas-armonicas}
\end{equation}
Division by three belongs to the reader of the three orbits. Its integrality
is verified on the compatible outputs of the record; it is not assumed for an
arbitrary element of \(\mathbb Z^{25}\).

The three signed blocks are
\begin{equation}
 U^{(r)}=(s_r,d_{r,1}-d_{r,3},d_{r,0}+d_{r,2},d_{r,0}-d_{r,2}).
 \label{eq:k-bloques-firmados}
\end{equation}
Integral substitution gives
\[
 (B_1,B_2,B_3)=(972,-1073,101),\qquad
 (s_1,s_2,s_3)=(2378,1406,2479),
\]
and therefore
\begin{align*}
 U^{(1)}&=(2378,-452,-668,-322),\\
 U^{(2)}&=(1406,998,-204,-28),\\
 U^{(3)}&=(2479,-551,-371,-997).
\end{align*}
Interleaving by columns recovers
\begin{equation}
 U=U_{\mathrm{term}}=(2378,1406,2479,-452,998,-551,
       -668,-204,-371,-322,-28,-997).
 \label{eq:k-u-firmado}
\end{equation}
The presence of negative values and components greater than 999 shows that
\(U\) is not a base-one-thousand digit word. Nor is it
\(U_6\Vert U_6\), a catalog concatenation of period six. The difference
concerns domain and preserved information, not merely values.

The complete representation and its return control are, respectively,
\[
\begin{aligned}
 X_{\mathrm{term}}
 &\xrightarrow{\mathcal R_{12}}
 \mathcal R_{12}(X_{\mathrm{term}})
 \xrightarrow{\mathcal E_{108}^{90,120}}
 C_{\mathrm{term}}
 \xrightarrow{\Pi_H}U_{\mathrm{term}}
 \xrightarrow{H_4^{-1}}K,\\
 K&\longrightarrow\mathscr F_{\mathrm{term}}^{(10)}
 \longrightarrow(b^{90},b^{120},Q_{\mathrm{TPK}})
 \longrightarrow U_{\mathrm{term}}.
\end{aligned}
\]
The Hadamard inversion develops \(U_{\mathrm{term}}\to K\), while the
incidence evaluation and harmonic reading develop the return
\(\mathscr F_{\mathrm{term}}^{(10)}\to U_{\mathrm{term}}\). Their equality is
a reversible concordance square; it does not turn a received table into the
generator of the signed state. The executable auxiliary serializes
\(C_{\mathrm{term}}\) in order to isolate and falsify this finite chart. The
initial equality is evaluated in the preceding passage through \(\Pi_H\) on
the coordinate \(C_{\mathrm{term}}\) already produced in the full state; the
auxiliary neither redefines the APP--TRIT--TPK generator nor factors the chart
through \(\mathsf{CT108RawProjection}\).

\begin{theorem}[Naturality of the signed reader]
\label{thm:k-naturalidad-lector-firmado}
For \(N'\geq N\geq108\), if
\[
 \tau_{N',N}:
 \mathcal X_{\mathrm{TPK}}^{[N'],\mathrm{enr}}
 \longrightarrow
 \mathcal X_{\mathrm{TPK}}^{[N],\mathrm{enr}}
\]
is the truncation that preserves the first \(N\) updates and their ledger,
then
\[
 \mathcal S_{12,N}\circ\tau_{N',N}
 =\mathcal S_{12,N'}.
\]
\end{theorem}

\begin{proof}
Each updater adds the next incidence without rewriting the one hundred eight
terminal windows already generated. Therefore \(\mathcal R_{12}\),
\(\mathcal E_{108}^{90,120}\), and \(\Pi_H\) read the same coordinates after
truncating any subsequent prolongation. The equality is one of mappings on
produced states; it does not result from preserving the output as a first
coordinate of the input.
\end{proof}

\subsection{Orbitwise Hadamard transform and integral reconstruction}
\label{subsec:k-hadamard}

In the twelve-position cycle, step three determines three orbits:
\[
 (1,4,7,10),\qquad(2,5,8,11),\qquad(3,6,9,12).
\]
The separation into four sectors of each orbit is what, once origin and
orientation have been fixed, admits the ninety-degree angular reading. The
transformation does not act on three contiguous tetrads in natural order.

Let
\[
 H_4=\begin{pmatrix}
 1&1&1&1\\1&1&-1&-1\\1&-1&1&-1\\1&-1&-1&1
 \end{pmatrix}.
\]
If \(P_{\mathrm{orb}}\) reorders according to the three stated orbits,
define
\[
 \mathcal H_{12}=P_{\mathrm{orb}}^{-1}
 (H_4\oplus H_4\oplus H_4)P_{\mathrm{orb}}.
\]

\begin{lemma}[Hadamard inversion and integrality]
\label{lem:k-hadamard}
One has
\[
 H_4^2=4I_4,\qquad H_4^{-1}=\frac14H_4,\qquad\det H_4=-16.
\]
For \(u\in\mathbb Z^4\), the preimage \(H_4^{-1}u\) is integral if and only
if \(H_4u\equiv0\pmod4\).
\end{lemma}

\begin{proof}
The rows are orthogonal and have squared norm four; the matrix is symmetric.
This gives \(H_4^2=4I_4\) and the inverse. Integral elimination, or expansion
of the determinant, gives the negative sign and the value \(-16\). The
integrality condition is exactly the divisibility by four of the four
coordinates of \(H_4u\).
\end{proof}

\begin{definition}[Integral dodecaphase transformation]
The subset of the integrality sublattice formed by the signed records and its integral coordinatization are
\[
 \mathcal L_H=\{u\in\ZZ^{12}:\mathcal H_{12}u\equiv0\pmod4\},
 \qquad \mathscr K(u)=\tfrac14\mathcal H_{12}u.
\]
The mapping \(\mathscr K:\mathcal L_H\to\ZZ^{12}\) is an isomorphism
of Abelian groups, with inverse \(k\mapsto\mathcal H_{12}k\).
The canonical record \(K\) is the image of the signed record generated
by the preceding operations, with phase origin and orientation fixed.
\end{definition}

\begin{proof}
The congruence defines exactly the integrality of \(\mathscr K(u)\).
The identity \(\mathcal H_{12}^2=4I\) proves that the two compositions are
inverse. In particular, every \(k\in\ZZ^{12}\) has preimage
\(\mathcal H_{12}k\in\mathcal L_H\).
\end{proof}

\begin{proposition}[Canonical integral record]
\label{prop:k-sello}
The inversion \(K^{(r)}=H_4U^{(r)}/4\) of the blocks in
\eqref{eq:k-bloques-firmados} produces
\begin{align*}
 K^{(1)}&=(234,729,621,794),\\
 K^{(2)}&=(543,659,58,146),\\
 K^{(3)}&=(140,824,914,601).
\end{align*}
In natural order, the dodecaphase register is
\begin{equation}
 K=(234,543,140,729,659,824,621,058,914,794,146,601).
 \label{eq:k-vector}
\end{equation}
It is the unique rational preimage of \(U\), and belongs to
\(\{0,\ldots,999\}^{12}\).
\end{proposition}

\begin{proof}
The three products before division are
\begin{align*}
 H_4U^{(1)}&=(936,2916,2484,3176),\\
 H_4U^{(2)}&=(2172,2636,232,584),\\
 H_4U^{(3)}&=(560,3296,3656,2404).
\end{align*}
All coordinates are divisible by four. Dividing and undoing the orbit
permutation gives \eqref{eq:k-vector}. Uniqueness follows from rational
invertibility, not from a search among candidates close to a physical value.
\end{proof}

Integrality can also be formulated in terms of lattices. The determinantal
divisors of \(H_4\) are \(1,2,4,16\): the entries have greatest common
divisor one; the minors of order two are even and one has modulus two; the
minors of order three are multiples of four and one has modulus four. Hence
\[
 \operatorname{SNF}(H_4)=\operatorname{diag}(1,2,2,4),
\]
and for the global transformation,
\[
 \operatorname{coker}\mathcal H_{12}
 \cong(\mathbb Z/2\mathbb Z)^6\oplus(\mathbb Z/4\mathbb Z)^3.
\]
The index of the image is \(16^3=4096\). Thus, not every signed integral
vector admits an integral preimage; the vector obtained satisfies the precise
congruences required by the inversion. The factor \(1/4\) is forced by the
matrix and does not have the status of a free coefficient.

\subsection{Cyclic differences and uniform component}
\label{subsec:k-transversal}

With cyclic indices modulo twelve, let
\[
 (D_3z)_m=z_m-z_{m+3},\qquad(D_4z)_m=z_m-z_{m+4},
 \qquad Q(z)=\sum_{m=1}^{12}z_m.
\]
The first operator compares sectors separated by three positions; the second,
by four. The channel names 90 and 120 refer here to those separations in the
oriented cycle, not to an identification with a physical operator from
another domain.

\begin{proposition}[Completeness of the transverse system]
\label{prop:k-transversal}
The operator
\[
 \mathcal D=(D_3,D_4,Q):\mathbb Z^{12}\longrightarrow\mathbb Z^{25}
\]
has rank twelve and nonzero invariant factors
\[
 \operatorname{diag}(1,\ldots,1,12),
\]
with eleven unit factors. In particular, its outputs determine a unique dodecaphase register.
\end{proposition}

\begin{proof}
If \(D_3z=D_4z=0\), the vector is constant along both steps.
Since \(3\) and \(4\) generate the cycle modulo twelve,
\(\ker D_3\cap\ker D_4=\langle\mathbf1\rangle\). The total charge
eliminates this final freedom because \(Q(\mathbf1)=12\).

For the integral assertion, the rows of \((D_3,D_4)\) are oriented
incidences of a connected graph. A spanning tree provides eleven unit
invariant factors. Every minor of order twelve must contain the row \(Q\).
Adding the first eleven columns to the last makes the latter vanish in the
incidence rows and equal twelve in row \(Q\). All maximal minors are divisible
by twelve, and the tree produces one of modulus exactly twelve. The final
factor is therefore twelve.
\end{proof}

Verification on \eqref{eq:k-vector} gives
\[
 D_3K=b^{90},\qquad D_4K=b^{120},\qquad Q(K)=6263.
\]
Formulas \eqref{eq:k-sumas-armonicas} and
\eqref{eq:k-bloques-firmados} in turn reconstruct \(U\) from those
differences. A triangle of exact representations of the same object is thus
closed: transverse record, signed vector, and dodecaphase register. Equality
of results does not turn their codomains into a single codomain; it proves
that the stated mappings recover the relevant information.

\subsection{Reading of one turn and periodic prolongation of the dodecaphase register}
\label{subsec:k-kappas}

Once \(K\) has been constructed, the base-one-thousand chart permits two
distinct operations. The first reads a single turn; the second periodically
prolongs the dodecaphase register. Let \(B=1000\) and define its concatenation integer
\begin{equation}
 N_K=\sum_{m=1}^{12}K_mB^{12-m}
 =234543140729659824621058914794146601.
 \label{eq:k-numerador}
\end{equation}
The leading zeros of a block, such as that of 058, are part of its three-digit
representation. They do not change its integral value, but they do permit the
twelve blocks to be concatenated unambiguously.

\begin{definition}[Finite and periodic readings of the dodecaphase register]
\label{def:k-kappas}
The reading of one turn is
\[
 \kappa_{36}=\sum_{m=1}^{12}K_mB^{-m}=\frac{N_K}{B^{12}}.
\]
The periodic prolongation is
\[
 K_j^{\mathrm{per}}=K_{1+((j-1)\bmod12)},\qquad
 \kappa_{\mathrm{per}}=\sum_{j\geq1}K_j^{\mathrm{per}}B^{-j}.
\]
\end{definition}

\begin{proposition}[Exact value and period of the prolonged dodecaphase register]
\label{prop:k-periodo}
One has
\begin{equation}
 \kappa_{\mathrm{per}}=\frac{N_K}{B^{12}-1},\qquad
 \kappa_{\mathrm{per}}-\kappa_{36}
 =\frac{N_K}{B^{12}(B^{12}-1)}
 =B^{-12}\kappa_{\mathrm{per}}.
 \label{eq:k-diferencia-kappas}
\end{equation}
The infinite word of the dodecaphase register has minimum period twelve in base
one thousand and minimum period thirty-six in base ten.
\end{proposition}

\begin{proof}
The sum of the first turn is \(N_KB^{-12}\); each new turn multiplies its
contribution by \(B^{-12}\). The geometric series yields \(N_K/(B^{12}-1)\),
and subtraction gives the second identity.

The twelve components of \(K\) are distinct. A proper period of the cycle
would identify two of them; therefore the minimum block period is twelve. The
decimal expansion is canonical, since the dodecaphase register has no tail of
999 blocks. If its minimum decimal period were \(d\), it would divide 36.
Grouping by threes would produce a block period \(d/\gcd(d,3)\). No proper
divisor of 36 gives twelve under that operation. Hence the minimum decimal
period is 36.
\end{proof}

Both readings are rational, but they are not the same rational number. The
first terminates after twelve blocks; the second repeats the turn. Moreover,
\(0<\kappa_{\mathrm{per}}-\kappa_{36}<B^{-12}\). Their proximity does not
permit one to be substituted for the other in an exact identity. In
particular, the finite closure of alpha uses \(\kappa_{36}\), whereas the sum
of the prolonged section uses \(\kappa_{\mathrm{per}}\).

\subsection{Reversibility of the rational constant and phase change}

The rational constant \(\kappa_{\mathrm{per}}\) completely determines
the record \(K\) when the base, length, and reading origin are preserved.
Its numerical function and its structural function are related by an
exact inversion.

\begin{proposition}[Integral recovery and cyclic transformation]
Let \(I(K)=\sum_{j=1}^{12}K_j1000^{12-j}\). Then
\[
 I(K)=(1000^{12}-1)\kappa_{\mathrm{per}}.
\]
The integral expansion of \(I(K)\) into twelve base-\(1000\) blocks
recovers every \(K_j\), including the leading zeros of each block.
If \(\rho(K)=(K_2,\ldots,K_{12},K_1)\), then
\begin{equation}
 \kappa(\rho K)=1000\kappa(K)-K_1.
 \label{eq:k-rotacion-racional}
\end{equation}
\end{proposition}
\begin{proof}
The first identity inverts the definition of the periodic evaluation.
Iterated Euclidean division by \(1000\), with twelve positions fixed,
recovers the vector. On the other hand,
\[
 I(\rho K)=1000I(K)-K_1(1000^{12}-1).
\]
Division by \(1000^{12}-1\) gives \eqref{eq:k-rotacion-racional}.
\end{proof}

The recovery therefore composes \(\kappa\mapsto K\mapsto U\)
and the already constructed transformations \(D_3K,D_4K,Q\). The rational
evaluation preserves the signed record and the cyclic differences that
depend on it. The total depth and the transport operators of a history
remain in their own coordinates: the domain of the preceding inversion
is the integral dodecaphase register.

Write \(\kappa_j=\kappa(\rho^jK)\), with indices modulo \(12\).
The same identity gives
\[
 K_{j+1}=1000\kappa_j-\kappa_{j+1},\qquad
 \sum_{j=0}^{11}\kappa_j=\frac{\sum_{j=1}^{12}K_j}{999}
 =\frac{6263}{999}.
\]
The phase acts exactly on the constant: its value is fixed
in the canonical reading, while its rotations follow an affine law.
The orbit sum is invariant under the choice of origin.

The subsequent incidence also has an expression in these
coordinates. The cubic threshold determines
\[
 B_K=\{j+1:1000\kappa_j-\kappa_{j+1}\ge729\}
 =\{4,6,9,10\}.
\]
The tetrad is recovered from the rational orbit because the latter
first recovers the integral components. The signed support associated
with \(\alpha\) also uses its regional records and the boundary
normalization. The relation between coordinate and incidence is thus
expressed through identifiable operations.

This is an operational definition of the function of \(K\): it provides
a reversible integral coordinatization of the oriented record, whose
periodic evaluation and incidence functions participate in subsequent
constructions. The value, transformation, and domain are preserved jointly.


\subsection{Periodicity and cokernel of the transport operator}
\label{subsec:k-clase-acarreo}

The periodicity just proved belongs to the dodecaphase register. It does not
imply periodicity of the regional coordinates, of the carries of the closure,
or of alpha. It is useful to specify this separation through a second, purely
integral quotient.

Let \(S\) be the cyclic shift of twelve coordinates, with orientation fixed,
and let \(b\geq2\) be an integer base. The periodic carry operator is
\(\partial_b=S-bI\). If a periodic identity is written
\[
 A=P+E-\Phi-K+\partial_bC,
\]
then the transformations
\[
 K\longmapsto K+\partial_bh,\qquad C\longmapsto C+h
\]
preserve its right-hand side. This invariance compares representatives of a
carry class; it is not an authorization to modify the dodecaphase register
already selected by the record.

\begin{proposition}[Quotient of the periodic carry]
\label{prop:k-cociente-acarreo}
With the preceding orientation,
\[
 \operatorname{coker}(S-bI)\cong\mathbb Z/(b^{12}-1)\mathbb Z.
\]
\end{proposition}

\begin{proof}
The cyclic module is presented as \(\mathbb Z[t]/(t^{12}-1)\). Imposing
\(S=bI\) is equivalent to imposing \(t=b\), and leaves the relation
\(b^{12}-1=0\). Evaluation of the coefficients at \(b\), with the chosen
concatenation order, represents the resulting class.
\end{proof}

The denominator \(B^{12}-1\) of \(\kappa_{\mathrm{per}}\) and this quotient
arise from the same cyclic length, but their objects remain distinct: one is a
rational evaluation and the other a quotient of a carry module. The finite
closure of alpha will be an oriented chain with a right boundary. Its
prolongation will transport a compatible remote boundary, not a carry
identified by decree between the first and thirteenth positions.

These results fix the internal inputs of the dodecaphase closure of
\(\alpha\). The dodecaphase register has been fixed before solving the alpha
recurrence; its Hadamard reconstruction is integral and unique; its
transverse differences recover the record; and its two rational evaluations
have separate domains and uses. None of these facts uses alpha as an input,
and none yet decides the arithmetic type of the output of the complete
closure.
